\documentclass[12pt]{article}

\usepackage[T1]{fontenc}
\usepackage[utf8]{inputenc}
\usepackage{lmodern}
\usepackage{microtype}
\usepackage[letterpaper,margin=1in]{geometry}
\usepackage{amssymb,amsthm,mathtools}
\usepackage{titlesec}
\usepackage[hidelinks]{hyperref}

\hypersetup{
  pdftitle={Nonquadraticity of Lambert Series Values},
  pdfsubject={Degree exclusion for rational and algebraic Lambert bases},
  pdfauthor={Costa Lambrinoudis}
}

\allowdisplaybreaks
\linespread{1.02}
\setlength{\parindent}{0pt}
\setlength{\parskip}{0.2\baselineskip}
\setlength{\skip\footins}{1.25\baselineskip}
\emergencystretch=2em
\titleformat{\section}{\raggedright\normalfont\Large\bfseries}{\thesection}{1em}{}
\titlespacing*{\section}{0pt}{3.5ex plus 1ex minus 0.2ex}{2.3ex plus 0.2ex}
\titleformat{\subsection}{\raggedright\normalfont\large\bfseries}{\thesubsection}{1em}{}
\titlespacing*{\subsection}{0pt}{3.25ex plus 1ex minus 0.2ex}{1.5ex plus 0.2ex}
\newtheorem{theorem}{Theorem}[section]
\newtheorem{lemma}[theorem]{Lemma}
\newtheorem{corollary}[theorem]{Corollary}
\newtheorem{proposition}[theorem]{Proposition}
\theoremstyle{remark}
\newtheorem{remark}[theorem]{Remark}

\newcommand{\Q}{\mathbb Q}
\newcommand{\Z}{\mathbb Z}
\newcommand{\C}{\mathbb C}
\newcommand{\R}{\mathbb R}
\newcommand{\pos}[1]{\left(#1\right)_+}
\DeclareMathOperator{\Res}{Res}
\DeclareMathOperator{\ord}{ord}
\begin{document}
\begin{center}
{\LARGE Nonquadraticity of Lambert Series Values}\\[1.5ex]
{\large Costa Lambrinoudis {\normalsize(with AI)}}\footnote{\textit{AI disclosure:} Generative AI was used extensively in the mathematical investigation, proof development, literature review, exposition, and Lean formalization, under the author's direction and review.}\\[0.5ex]
{7 October 2026}
\end{center}

\begin{quote}\small
For every integer $q\geq2$, the Lambert value
$F(q)=\sum_{n\geq1}(q^n-1)^{-1}$ satisfies no nonzero polynomial of degree at
most two over $\Q$. In particular, the Erd\H{o}s--Borwein constant $F(2)$ is
neither rational nor quadratic irrational. More generally, for real algebraic
$q>1$, we exclude every nonzero polynomial relation of degree at most $e$ over
$\Q(q)$ whenever $e\lambda\log M(q)<\log q$, where $M(q)$ is the Mahler measure
of the base and $\lambda=(42963+428220/\pi^2)/202140$.
This gives nonquadraticity over $\Q(q)$ for Pisot and Salem bases,
$F(\sqrt m)\notin\Q(\sqrt m)$ for every integer $m\geq2$, and irrationality
for rational bases $q=a/b>1$ with $\log b/\log a<0.5728$.
The proof combines positive moment determinants, local denominator bounds,
and the product formula.
\end{quote}

\section{Introduction}

The Lambert series
\[
F(q)=\sum_{n=1}^{\infty}\frac1{q^n-1},\qquad q>1,
\]
converges absolutely. Erd\H{o}s proved that $F(q)$ is irrational at every integer
base $q\geq2$ in 1948~\cite{erdos1948}. We prove nonquadraticity at every
integer base, including the Erd\H{o}s--Borwein constant $F(2)$.

\begin{theorem}[Integer bases]\label{thm:integer}
For every integer $q\geq2$, the numbers $1,F(q),F(q)^2$ are linearly
independent over $\Q$.
\end{theorem}

We obtain this result from a criterion over $\Q(q)$ for arbitrary real
algebraic bases $q>1$. The criterion compares the degree of a proposed relation
with $\log q/\log M(q)$, where $M(q)$ is the Mahler measure of the base.
It also applies to Pisot, Salem, and radical bases.

For rational bases $q=a/b>1$, the same theorem improves the known bounds for
Erd\H{o}s Problem 1049,\footnote{\url{https://www.erdosproblems.com/1049}.}
enlarging the irrationality region from $\log b/\log a<0.405683\ldots$ to
$\log b/\log a<0.5728$. See Erd\H{o}s~\cite[p.~102]{erdos} for the problem.
The earlier bound of
Bundschuh--V\"a\"an\"anen~\cite[Theorem~2]{bundschuh-vaananen} was
$1/2-1/\pi^2=0.398678\ldots$; Cook's 2026 refinement of Zudilin's
permutation method~\cite{cook,zudilin} gives the former threshold.

Positive moment determinants provide nonvanishing and decay at the distinguished
real embedding. Local denominator bounds produce integral polynomials, while
uniform complex estimates control their values at the other conjugates.
The product formula then excludes algebraic relations of low degree.
The method belongs to the same family as Fauzan's 2026 preprint on
$\zeta(5)$~\cite{fauzan}: positive moment determinants with local denominator
control, with cyclotomic arithmetic as the $q$-specific ingredient here.

\subsection{The general degree criterion}

For a real algebraic number $q$, put $K=\Q(q)$ and let
\[
M(q)=|a_g|\prod_{j=1}^{g}\max(1,|q_j|)
\]
be the Mahler measure of its primitive irreducible integer polynomial
$a_g\prod_j(X-q_j)$. Define
\begin{equation}\label{eq:constants}
C=42963+\frac{428220}{\pi^2},\qquad S=202140,\qquad
\lambda=\frac CS=0.4271829289\ldots.
\end{equation}

\begin{theorem}\label{thm:algebraic}
Let $q>1$ be real algebraic and $e\geq0$ an integer. If
\begin{equation}\label{eq:algmargin}
e\lambda\log M(q)<\log q,
\end{equation}
then $P(F(q))\ne0$ for every nonzero $P\in K[X]$ of degree at most $e$.
\end{theorem}

\begin{proof}[Proof of Theorem~\ref{thm:integer}]
For an integer $q\geq2$, $K=\Q$ and $M(q)=q$. Since $2\lambda<1$,
Theorem~\ref{thm:algebraic} applies with $e=2$.
\end{proof}

\begin{corollary}\label{cor:conjugates}
If $q>1$ is an algebraic integer whose other conjugates have modulus at most one, then $1,F(q),F(q)^2$ are linearly independent over $\Q(q)$.
\end{corollary}
\begin{proof}
Here $M(q)=q$, and $2C<S$. Apply Theorem~\ref{thm:algebraic} with $e=2$.
\end{proof}
In particular, this applies to every Pisot or Salem base.

\begin{corollary}[Rational bases]\label{cor:rational}
Let $a>b>0$ be integers and let $e\geq0$ be an integer. If
\[
\frac{\log b}{\log a}<1-e\lambda,
\]
then no nonzero polynomial in $\Q[X]$ of degree at most $e$ vanishes at $F(a/b)$.
In particular, $F(a/b)$ is irrational whenever $\log b/\log a<0.5728$.
\end{corollary}
\begin{proof}
For coprime $a,b$, $M(a/b)=a$, so the assertion follows from
Theorem~\ref{thm:algebraic}. In general $M(a/b)\leq a$, which gives the same
sufficient condition. Finally $1-\lambda=0.5728170710\ldots>0.5728$.
\end{proof}

\begin{corollary}\label{cor:radical}
Suppose $q>1$ and $q^n=m$ for positive integers $n,m$. Put $g=[\Q(q):\Q]$.
If $e\geq0$ is an integer and $eg\lambda<1$, no nonzero polynomial over $\Q(q)$ of degree at most $e$ vanishes at $F(q)$.
In particular,
\[
F(\sqrt m)\notin\Q(\sqrt m)\qquad(m\in\Z,\ m\geq2).
\]
\end{corollary}
\begin{proof}
The number $q$ is integral and every conjugate has modulus $q$, since its $n$th power is $m$.
Hence $\log M(q)=g\log q$. For the last assertion take $e=1$ and use $g\leq2$ and $2\lambda<1$.
\end{proof}

\section{Positive determinants and their normalization}\label{sec:determinants}

For an indeterminate base $U$, define a linear functional by
\begin{equation}\label{eq:functional}
\mu_X(t^r)=\frac1{U^{r+1}-1},\qquad
\mu_X\!\left(\frac1{U^j-t}\right)=X-S_j(U),\qquad
S_j(U)=\sum_{v=1}^{j}\frac1{U^v-1}.
\end{equation}
Partial fractions extend it to rational functions with distinct poles among the powers of $U$.
For integers $h\geq1$ and $k\geq0$, put
\begin{equation}\label{eq:delta}
B_{h,k}(t)=\prod_{j=k}^{k+h-1}(U^j-t),\qquad
\Delta_{h,k}(U,X)=\det_{0\leq i,j<h}
\mu_X\!\left(\frac{t^{k+i+j}}{B_{h,k}(t)}\right).
\end{equation}
Each entry is affine in $X$, so $\deg_X\Delta_{h,k}\leq h$.
The parameter $k$ shifts both the poles and the moments.

For $q>1$, let $d\nu_q=\sum_{v\geq1}q^{-v}\delta_{q^{-v}}$. Direct summation gives
\[
\int t^r\,d\nu_q(t)=\frac1{q^{r+1}-1},\qquad
\int\frac{d\nu_q(t)}{q^j-t}=F(q)-S_j(q).
\]
On its support $0<t\leq q^{-1}$, every factor $q^j-t$ is positive, so the
weight $t^k/B_{h,k}(t)$ is strictly positive. Evaluation of~\eqref{eq:delta} at
$(q,F(q))$ is therefore the Gram determinant of $1,t,\ldots,t^{h-1}$ for a
positive measure with infinite support. The matrix is positive definite.

Define
\begin{equation}\label{eq:normalization}
G_{h,k}=k\left(hk+\frac{h(h-1)}2\right),\qquad
Q_{h,k}(U,X)=U^{-G_{h,k}}\Delta_{h,k}(U,X).
\end{equation}
This normalization makes every coefficient bounded at infinity.

\begin{proposition}\label{prop:analytic}
For $q>1$,
\begin{equation}\label{eq:decay}
Q_{h,k}(q,F(q))>0,\qquad
\log Q_{h,k}(q,F(q))=-\mathcal T_{h,k}\log q+O_q(h),
\end{equation}
where
\begin{equation}\label{eq:T}
\mathcal T_{h,k}=G_{h,k}+h\sum_{j=k}^{k+h-1}j+
\frac{h(h+1)(2h+1)}6+\frac{k h(h+1)}2.
\end{equation}
For every fixed $c_0>0$, $\rho>1$, and $Z\geq0$, if $h,k\leq c_0N$, then uniformly
for $|u|=\rho$ and $|z|\leq Z$,
\begin{equation}\label{eq:complexQ}
|Q_{h,k}(u,z)|\leq\exp\bigl(O_{\rho,Z,c_0}(N\log(N+1))\bigr).
\end{equation}
Every coefficient of $Q_{h,k}$ has degree at infinity at most zero.
\end{proposition}
\begin{proof}
Put $\gamma_q=\prod_{v\geq1}(1-q^{-v})>0$. On the support of $\nu_q$,
\[
q^{-\sum_{j=k}^{k+h-1}j}\leq B_{h,k}(t)^{-1}
\leq\gamma_q^{-1}q^{-\sum_{j=k}^{k+h-1}j}.
\]
Comparison of positive Gram forms bounds $\Delta(q,F(q))$ between the
corresponding $h$th powers of these bounds times
\[
\det_{0\leq i,j<h}\frac1{q^{k+i+j+1}-1}.
\]
The Cauchy determinant formula factors this expression into its leading power
\[q^{-h(h+1)(2h+1)/6-kh(h+1)/2}\]
and products of factors $1-q^{-v}$.
Their logarithms contribute $O_q(h)$, uniformly in $k\geq0$.
Multiplication by $q^{-G_{h,k}}$ proves~\eqref{eq:decay}.

For the remaining assertions put
\[
\mathcal C_{i,j}(U,X)=U^{ij}(X-S_j(U))-
\sum_{v=1}^{i}\frac{U^{j(i-v)}}{U^v-1},\qquad
V_{h,k}(U)=\prod_{0\leq i<j<h}(U^{k+j}-U^{k+i}).
\]
Partial fractions, followed by the change from the simple-pole basis to
$1,t,\ldots,t^{h-1}$, give
\begin{equation}\label{eq:augmented}
\det_{0\leq i,j<h}\mathcal C_{k+i,k+j}(U,X)
=(-1)^{\binom h2}V_{h,k}(U)\Delta_{h,k}(U,X).
\end{equation}
The entry $\mathcal C_{i,j}$ has degree at infinity at most $ij$. Rearrangement shows
that the largest permutation degree in the determinant is
$\sum_{i<h}(k+i)^2$. Subtracting
$\deg V_{h,k}=\sum_{i<h}i(k+i)$ leaves exactly $G_{h,k}$.
This proves the coefficientwise infinity bound.

For $\rho=|u|>1$,
\[
|\mathcal C_{i,j}(u,z)|\leq\rho^{ij}\left(|z|+\frac{i+j}{\rho-1}\right),\qquad
|V_{h,k}(u)|\geq\gamma_\rho^h\rho^{\sum_{i<h}i(k+i)}.
\]
Apply the same rearrangement to each permutation term in~\eqref{eq:augmented}
and cancel $\rho^{G_{h,k}}$. The remaining factorial and entry factors have
logarithm $O_{\rho,Z,c_0}(N\log(N+1))$, proving~\eqref{eq:complexQ}.
\end{proof}

\section{Local denominator bounds}\label{sec:arithmetic}

Identity~\eqref{eq:augmented} shows that coefficient denominators are supported
at $U=0$ and the cyclotomic polynomials $\Phi_n(U)$. We bound their orders
before specializing the base.

\subsection{The zero place}

Work in the Laurent series field at $U=0$, and put
\[
P(t)=\prod_{v=k}^{k+h-1}(t-U^v),\qquad x_j=U^j.
\]
At a pole $x_j$, write $P_j(t)=P(t)/(t-x_j)$. The finite part of $p(t)/P(t)$ there is
\[
\frac{p'(x_j)}{P_j(x_j)}-
\frac{p(x_j)P_j'(x_j)}{P_j(x_j)^2}.
\]
For an integer $R>k+h$, put
\[
Z_R=X+\sum_{v=1}^{R}\frac{U^v}{1-U^v}.
\]
The functional obtained by summing $-x_j$ times these finite parts for $0\leq j<R$, and subtracting
$Z_R\sum_{k\leq j<k+h}p(x_j)/P_j(x_j)$, converges to
$\mu_X(p/P)$. Indeed, polynomial division reduces this statement to
\[
-\sum_{j=0}^{R-1}U^{j(r+1)}\longrightarrow\frac1{U^{r+1}-1}
\]
and the simple-pole identity obtained by separating $j<v$ and $j>v$ in
$\sum_{j<R,\,j\ne v}U^j/(U^j-U^v)$. The correction $Z_R$ gives the limit
$S_v(U)-X$ at the pole $t=U^v$.

Write this finite functional as
$\sum_{j<R}\{W_jp(x_j)+\Gamma_jx_jp'(x_j)\}$. Its weights are
\[
(W_j,\Gamma_j)=
\begin{cases}
\displaystyle\left(-\frac{x_j}{P(x_j)},0\right),&j\notin[k,k+h),\\[1.2ex]
\displaystyle\left(\frac{x_jP_j'(x_j)/P_j(x_j)-Z_R}{P_j(x_j)},
-\frac1{P_j(x_j)}\right),&j\in[k,k+h).
\end{cases}
\]
Since $\ord(U^j-U^v)=\min(j,v)$ for $j\ne v$, and both $Z_R$ and
$x_jP_j'(x_j)/P_j(x_j)$ have nonnegative order, each weight has order at least
\[
\omega_j=j-\sum_{v=k}^{k+h-1}\min(j,v).
\]
The resulting finite moment matrix has entries
$\sum_{j<R}(W_j+(k+i+r)\Gamma_j)x_j^{k+i+r}$. At each pole this is the product of the two-column block
$\bigl(x_j^i,\,i x_j^i\bigr)$ and the two-row block
$\bigl((W_j+(k+r)\Gamma_j)x_j^{k+r},\,\Gamma_jx_j^{k+r}\bigr)^{\mathsf T}$.
Outside the pole window only the evaluation column remains. The sign change
$P=(-1)^hB_{h,k}$ does not affect orders.

\begin{lemma}\label{lem:zero}
Put
\[
w(j,r)=(k+1)j-\sum_{v=k}^{k+h-1}\min(j,v)+2(h-1-r)j,
\qquad 0\leq r<h.
\]
A lower bound for $\ord_{U=0}\Delta$ is any lower bound for $\sum_{r<h}w(j_r,r)$ over admissible selections. It is sufficient to impose the relaxed capacity conditions
\[
j_r\geq r\quad(r<k),\qquad
j_r\geq k+\left\lfloor\frac{r-k}{2}\right\rfloor\quad(r\geq k).
\]
The same lower bound holds for every coefficient in $X$.
\end{lemma}
\begin{proof}
In each finite block, Cauchy--Binet expands the determinant over distinct columns. Before the pole window there is only one slot per node; allowing two slots per node thereafter can only enlarge the set of selections. This gives the displayed inequalities. Order the selected node indices increasingly as $j_0,\ldots,j_{h-1}$. In each of the two minors, every permutation of the exponents $0,\ldots,h-1$ has cost at least $\sum_r(h-1-r)j_r$, by rearrangement. The two minors therefore contribute at least twice this sum. The shift contributes $k\sum_rj_r$, and the weight contributes the remaining terms of $w$.

A sum has valuation at least the minimum valuation of its summands. The bounds therefore hold for every finite block and pass to the Laurent-series limit. To obtain coefficientwise bounds, evaluate $X$ at $h+1$ distinct rational constants, each of valuation zero, and interpolate: the inverse Vandermonde has constant entries and introduces no negative order. This gives the coefficientwise bound.
\end{proof}

\subsection{Cyclotomic places}

Write $x_+=\max(x,0)$ and
\[
v_h(n)=\sum_{j\geq1}(h-jn)_+,\qquad
\tau_{h,k}(n)=(k+h-\max(k,n))_+.
\]
The first quantity is the order of $V_{h,k}(U)$ at $\Phi_n$; the second counts
the indices in $[k,k+h)$ that are at least $n$.

\begin{lemma}\label{lem:cyclotomic}
Every coefficient of $\Delta_{h,k}$ has cyclotomic pole order at most
\begin{equation}\label{eq:cyclobound}
h+v_h(n)-\max\{v_h(2n),\ h-2\tau_{h,k}(n),\ 0\}.
\end{equation}
\end{lemma}
\begin{proof}
At a primitive $n$th root $\zeta$, use the local coordinate $U=\zeta e^t$.
The square matrix
\[
\mathcal M(t)=\bigl(t\mathcal C_{i,j}(U,X)\bigr)_{k\leq i,j<k+h}
\]
is regular at $t=0$, and~\eqref{eq:augmented} gives
$\det\mathcal M=\pm t^hV_{h,k}\Delta_{h,k}$.
It remains to bound the order of $\det\mathcal M$ in two ways.

First,
\[
\frac{t\mathcal C_{i,j}(U,X)}{U^{ij}}
=tX-\sum_{v=1}^{j}\frac{t}{U^v-1}
-\sum_{v=1}^{i}\frac{tU^{-jv}}{U^v-1}.
\]
Its constant term is $-(H_{\lfloor i/n\rfloor}+H_{\lfloor j/n\rfloor})/n$,
where $H_m=\sum_{v=1}^m1/v$.
For each $\ell\geq1$, the coefficient of $t^\ell$ in
$tU^{-jv}/(U^v-1)$ is, on each residue class of $v$ modulo $n$, a polynomial
in $v$ of degree at most $\ell-1$. For $n\mid v$, expand
$v^{-1}(vt/(e^{vt}-1))e^{-jvt}$; otherwise the denominator has nonzero
constant term. Summation over $1\leq v\leq i$ raises the degree by at most one.
Thus, after adding the row harmonic term, each coefficient of order $\ell$
is polynomial in $i$ of degree at most $\ell$ on each residue class.
Multiplication by $U^{ij}$ preserves this bound by convolution.

Consequently the coefficient vectors of order $\ell$ of every column of $\mathcal M$
lie in the span of
\[
i^a\mathbf1_{i\equiv v\pmod n},\quad
H_{\lfloor i/n\rfloor}i^a\mathbf1_{i\equiv v\pmod n},
\qquad 0\leq a\leq\ell,\quad0\leq v<n.
\]
These nested spaces have dimension at most $2n(\ell+1)$.
In a nonzero term of the determinant's Taylor expansion, at most $2n\ell$
chosen coefficient vectors can have order below $\ell$.
Its total order is therefore at least
$\sum_{\ell\geq1}(h-2n\ell)_+=v_h(2n)$.

Second, the constant matrix has entries
\[
\mathcal M(0)_{i,j}=-\frac{\zeta^{ij}}n
\left(H_{\lfloor i/n\rfloor}+H_{\lfloor j/n\rfloor}\right).
\]
The first summand is supported on $\tau_{h,k}(n)$ rows and the second on
that many columns. Its rank is at most $2\tau_{h,k}(n)$, so its corank is
at least $(h-2\tau_{h,k}(n))_+$. Constant invertible row operations make
that many rows divisible by $t$, giving the same lower bound on $\ord\det\mathcal M$.
Combining the two bounds and subtracting $h+v_h(n)$ proves the lemma.
All operations retain $X$ as a polynomial parameter, so the assertion holds
coefficientwise.
\end{proof}

\subsection{Integer descent}

By~\eqref{eq:augmented}, every coefficient of $Q_{h,k}$ lies in
\[
\mathcal R=\Z[U^{\pm1},(U^m-1)^{-1}:m\geq1],
\]
since the determinant identity divides only by powers of $U$ and differences
of powers. Let $J(U)$ be a monic integer polynomial such that every
coefficient of $JQ_{h,k}$ has nonnegative order at $U=0$ and at each
cyclotomic factor. These are all possible finite poles, so $JQ_{h,k}\in\Q[U,X]$.
For a coefficient of $Q_{h,k}$ written as $A(U)/B(U)$ with $A\in\Z[U]$ and $B\in\Z[U]$
monic, polynomiality says that $B$ divides $JA$ over $\Q[U]$.
Gauss's lemma then gives the quotient in $\Z[U]$. Hence
\begin{equation}\label{eq:integerpoly}
\mathcal P(U,X)=J(U)Q_{h,k}(U,X)\in\Z[U,X],\qquad
\deg_U\mathcal P\leq\deg J,\quad\deg_X\mathcal P\leq h.
\end{equation}
The first degree bound follows from Proposition~\ref{prop:analytic}.

\section{The arithmetic and analytic scales}\label{sec:scales}

Take $h=60N$ and $k=3N$, with $N\geq1$, and write $Q_N=Q_{60N,3N}$.
Equations~\eqref{eq:normalization} and~\eqref{eq:T} give
\[
G_N=5940N^3-90N^2,\qquad \mathcal T_{60N,3N}/N^3\longrightarrow S.
\]
We clear the local denominators by
\[
J_N(U)=U^{a_N}\prod_{n\geq1}\Phi_n(U)^{c_N(n)}.
\]

For each slot $r$, the cost in Lemma~\ref{lem:zero} is minimized under its
relaxed lower bound on $j_r$ by
\[
j_r=\begin{cases}r,&r<3N,\\
\max\{\lfloor(r+3N)/2\rfloor,\,2r-60N\},&3N\leq r<60N.
\end{cases}
\]
Inside the pole window, twice the cost is
$j(j+120N-1-4r)+(3N)(3N-1)$.
The displayed choice minimizes this integral quadratic subject to the slot bound;
the exterior affine pieces give no smaller value. Minimizing separately in each
slot gives a lower bound for every admissible selection.
The three ranges are $0\leq r<3N$, $3N\leq r<41N$, and
$41N\leq r<60N$. In the middle range write $r=3N+2t$ or $3N+2t+1$,
so that $j_r=3N+t$ for $0\leq t<19N$. In the last range $j_r=2r-60N$.
The standard sums of integers and squares give
\begin{align*}
\sum_{r=0}^{3N-1}w(j_r,r)
 &=\frac{531}{2}N^3-90N^2+\frac12N,\\
\sum_{r=3N}^{41N-1}w(j_r,r)
 &=\frac{20216}{3}N^3-361N^2+\frac{19}{3}N,\\
\sum_{r=41N}^{60N-1}w(j_r,r)
 &=-\frac{102163}{6}N^3+361N^2+\frac{19}{6}N.
\end{align*}
Adding these expressions gives
\[
-\sum_{r<60N}w(j_r,r)=10023N^3+90N^2-10N.
\]
Adding $G_N$ yields the clearing exponent
\begin{equation}\label{eq:algzero}
a_N=15963N^3-10N.
\end{equation}

For cyclotomic places with $2n\leq h$, use the filtration term of
Lemma~\ref{lem:cyclotomic}, giving $h+v_h(n)-v_h(2n)$.
For $2n>h$, we have $v_h(n)=(h-n)_+$, $v_h(2n)=0$, and
$\tau_{h,k}(n)=(63N-n)_+$, since $n>30N>k$.
The same lemma now gives
\[
(h-n)_++\min\{h,2(h+k-n)_+\}.
\]
These two bounds are nonnegative and vanish for $n\geq63N$.
To combine them, put
\[
E_h(n)=3(h-n)_+-(h-2n)_++\sum_{j\geq1}(h-(2j+1)n)_+.
\]
For $2n\leq h$, separating the odd terms of $v_h(n)$ shows that
$E_h(n)=h+v_h(n)-v_h(2n)$. For $2n>h$, it reduces to $3(h-n)_+$.
The identity $\min(h,2x_+)=2x_+-2(x-h/2)_+$ therefore gives the common formula
\begin{equation}\label{eq:alge}
\begin{split}
c_N(n)=E_{60N}(n)&+2(63N-n)_+-2(33N-n)_+\\
&-2(60N-n)_++2(30N-n)_+.
\end{split}
\end{equation}

Let $d_N=\deg J_N=a_N+\sum_n\varphi(n)c_N(n)$ and $u_N=\sum_n c_N(n)$.
The summatory totient estimate
$\sum_{n\leq x}\varphi(n)\sim3x^2/\pi^2$~\cite[27.11.6]{dlmf}, by partial summation, gives
\[
\frac1{N^3}\sum_n\varphi(n)(KN-mn)_+
\longrightarrow\frac{K^3}{\pi^2m^2}.
\]
The sums are bounded uniformly by a constant times $m^{-2}$, so this limit
may be summed over $m$. Using
$\sum_{j\geq1}(2j+1)^{-2}=\pi^2/8-1$, we obtain
\[
\frac1{N^3}\sum_n\varphi(n)E_{60N}(n)
\longrightarrow27000+\frac{378000}{\pi^2}.
\]
The four translated tents contribute
$2(63^3-33^3-60^3+30^3)/\pi^2=50220/\pi^2$.
Together with~\eqref{eq:algzero}, this proves
\begin{equation}\label{eq:alglimits}
\frac{d_N}{N^3}\longrightarrow15963+27000+\frac{378000+50220}{\pi^2}=C,
\qquad u_N=o(N^3).
\end{equation}
For the last assertion, use $c_N(n)\leq60N+v_{60N}(n)$ on $n<63N$ and
$v_{60N}(n)\leq(60N)^2/(2n)$, giving $u_N=O(N^2\log(N+1))$.

By~\eqref{eq:integerpoly},
$\mathcal P_N(U,X)=J_N(U)Q_N(U,X)\in\Z[U,X]$, with degrees at most $d_N$ in
$U$ and $60N$ in $X$. For fixed $q>1$, M\"obius inversion gives
$\log\Phi_n(q)=\varphi(n)\log q+O_q(1)$, uniformly in $n$.
Consequently $\log J_N(q)=d_N\log q+O_q(u_N)$, and Proposition~\ref{prop:analytic}
gives
\begin{equation}\label{eq:algsmall}
\mathcal P_N(q,F(q))>0,\qquad
\frac{\log\mathcal P_N(q,F(q))}{N^3}\longrightarrow(C-S)\log q.
\end{equation}

\section{The product-formula argument}\label{sec:algebraic}

\begin{lemma}\label{lem:disk}
For fixed $u\in\C$, $Z\geq0$, and $\epsilon>0$, eventually, uniformly for $|z|\leq Z$,
\[
|\mathcal P_N(u,z)|\leq
\exp\bigl((C\log\max(1,|u|)+\epsilon)N^3\bigr).
\]
\end{lemma}
\begin{proof}
For $\rho>1$, M\"obius inversion and the power-series expansion of $\log(1-w)$ give
\[
\bigl|\log|\Phi_n(u)|-\varphi(n)\log\rho\bigr|\leq c_\rho
\qquad(|u|=\rho),
\]
with $c_\rho$ independent of $n$ and the argument of $u$. Proposition~\ref{prop:analytic} and~\eqref{eq:alglimits} imply the bound $\exp((C\log\rho+o(1))N^3)$ on that circle. Maximum modulus extends it to $|u|\leq\rho$. Given $u$ and $\epsilon$, choose $\rho>\max(1,|u|)$ sufficiently close to that maximum, then take $N$ large. Only the cleared polynomial is evaluated inside the disk, so zero and roots of unity cause no singularity.
\end{proof}

For a number field $E$, use the unnormalized logarithmic height
\[
H_E(x)=\sum_{\sigma:E\hookrightarrow\C}\log\max(1,|\sigma(x)|)
+\sum_{v\ \text{finite}}\log\max(1,|x|_v),
\]
where finite-place absolute values include the local degrees, so the product formula has no additional weights. Complex embeddings are counted individually.

\begin{proposition}\label{prop:height}
If a real number field $E$ contains both $q>1$ and $F(q)$, then
\[
S\log q\leq CH_E(q).
\]
\end{proposition}
\begin{proof}
Put $\alpha=F(q)$ and $x_N=\mathcal P_N(q,\alpha)\ne0$. At each finite place, integrality of the polynomial coefficients gives
\[
\log|x_N|_v\leq d_N\log\max(1,|q|_v)
+60N\log\max(1,|\alpha|_v).
\]
The total $\alpha$ contribution is $O(N)$. At the distinguished real embedding use~\eqref{eq:algsmall}; at every other embedding use Lemma~\ref{lem:disk}. Sum the estimates, divide by $N^3$, and apply the product formula. Since the number of embeddings is fixed, their $\epsilon$ errors can be sent to zero, giving
$0\leq CH_E(q)-S\log q$.
\end{proof}

\begin{proof}[Proof of Theorem~\ref{thm:algebraic}]
The case $e=0$ is immediate. If a nonzero polynomial of degree at most $e\geq1$ over $K=\Q(q)$ vanished at $\alpha=F(q)$, then $E=K(\alpha)$ would have relative degree $r\leq e$. The height--Mahler identity~\cite[Proposition~1.6.6]{bombieri-gubler} and the tower formula give
\[
H_E(q)=r\log M(q).
\]
Proposition~\ref{prop:height} would imply
$S\log q\leq rC\log M(q)$, contradicting~\eqref{eq:algmargin}. For nonintegral bases, the finite-place contribution is accounted for by the leading coefficient in $M(q)$.
\end{proof}

\begin{remark}
Transporting the local bounds through six exact bordered presentations of the
same determinant improves the rational-base irrationality threshold to
$0.573290$; this refinement is formalized in the accompanying Lean development.
Numerical parameter searches suggest that this family of estimates saturates
near that value.
\end{remark}

\begin{thebibliography}{9}
\bibitem{erdos1948}
P.~Erd\H{o}s, \emph{On arithmetical properties of Lambert series},
Journal of the Indian Mathematical Society (N.S.) \textbf{12} (1948), 63--66.
\url{https://users.renyi.hu/~p_erdos/1948-04.pdf}.
\bibitem{erdos}
P.~Erd\H{o}s, \emph{On the irrationality of certain series: problems and results},
in A.~Baker (ed.), New Advances in Transcendence Theory,
Cambridge University Press, Cambridge, 1988, pp.~102--109.
\url{https://doi.org/10.1017/CBO9780511897184.009}.
\bibitem{bundschuh-vaananen}
P.~Bundschuh and K.~V\"a\"an\"anen, \emph{Arithmetical investigations of a certain infinite product},
Compositio Mathematica \textbf{91} (1994), no.~2, 175--199.
\url{https://www.numdam.org/item/CM_1994__91_2_175_0/}.
\bibitem{zudilin}
W.~Zudilin, \emph{Heine's basic transform and a permutation group for $q$-harmonic series},
Acta Arithmetica \textbf{111} (2004), no.~2, 153--164.
\url{https://doi.org/10.4064/aa111-2-4}.
\bibitem{cook}
W.~Cook, \emph{Geometric Moments and Rational Lambert Values: Proofs and Further Results},
online manuscript, September 2026.
\href{https://wcook04.github.io/plectis/maths/papers/erdos1049-rational-base-lambert-reasoning-surface.html}{Manuscript and source}, accessed 7 October 2026.
\bibitem{fauzan}
A.~Fauzan, \emph{$\zeta(5)$ is irrational}, preprint, 17 September 2026.
\url{https://doi.org/10.5281/zenodo.22826419}.
\bibitem{bombieri-gubler}
E.~Bombieri and W.~Gubler, \emph{Heights in Diophantine Geometry},
New Mathematical Monographs \textbf{4}, Cambridge University Press, 2006,
Proposition~1.6.6.
\url{https://doi.org/10.1017/CBO9780511542879}.
\bibitem{dlmf}
NIST Digital Library of Mathematical Functions, \S27.11, equation~27.11.6.
\url{https://dlmf.nist.gov/27.11.E6}.
\end{thebibliography}
\end{document}
